%%%%%%%%%%%%%%%%%%%%%%%%%%%%%%%%%%%%%%%%%
% The Legrand Orange Book
% LaTeX Template
% Version 2.4 (26/09/2018)
%
% This template was downloaded from:
% http://www.LaTeXTemplates.com
%
% Original author:
% Mathias Legrand (legrand.mathias@gmail.com) with modifications by:
% Vel (vel@latextemplates.com)
%
% License:
% CC BY-NC-SA 3.0 (http://creativecommons.org/licenses/by-nc-sa/3.0/)
%
% Compiling this template:
% This template uses biber for its bibliography and makeindex for its index.
% When you first open the template, compile it from the command line with the 
% commands below to make sure your LaTeX distribution is configured correctly:
%
% 1) pdflatex main
% 2) makeindex main.idx -s StyleInd.ist
% 3) biber main
% 4) pdflatex main x 2
%
% After this, when you wish to update the bibliography/index use the appropriate
% command above and make sure to compile with pdflatex several times 
% afterwards to propagate your changes to the document.
%
% This template also uses a number of packages which may need to be
% updated to the newest versions for the template to compile. It is strongly
% recommended you update your LaTeX distribution if you have any
% compilation errors.
%
% Important note:
% Chapter heading images should have a 2:1 width:height ratio,
% e.g. 920px width and 460px height.
%
%%%%%%%%%%%%%%%%%%%%%%%%%%%%%%%%%%%%%%%%%

%----------------------------------------------------------------------------------------
%	PACKAGES AND OTHER DOCUMENT CONFIGURATIONS
%----------------------------------------------------------------------------------------

\documentclass[11pt,fleqn]{book} % Default font size and left-justified equations

%%%%%%%%%%%%%%%%%%%%%%%%%%%%%%%%%%%%%%%%%
% The Legrand Orange Book
% Structural Definitions File
% Version 2.1 (26/09/2018)
%
% Original author:
% Mathias Legrand (legrand.mathias@gmail.com) with modifications by:
% Vel (vel@latextemplates.com)
% 
% This file was downloaded from:
% http://www.LaTeXTemplates.com
%
% License:
% CC BY-NC-SA 3.0 (http://creativecommons.org/licenses/by-nc-sa/3.0/)
%
%%%%%%%%%%%%%%%%%%%%%%%%%%%%%%%%%%%%%%%%%

%----------------------------------------------------------------------------------------
%	VARIOUS REQUIRED PACKAGES AND CONFIGURATIONS
%----------------------------------------------------------------------------------------

\usepackage{graphicx} % Required for including pictures
\graphicspath{{Pictures/}} % Specifies the directory where pictures are stored

\usepackage{lipsum} % Inserts dummy text

\usepackage{tikz} % Required for drawing custom shapes

\usepackage[english]{babel} % English language/hyphenation

\usepackage{enumitem} % Customize lists
\setlist{nolistsep} % Reduce spacing between bullet points and numbered lists

\usepackage{booktabs} % Required for nicer horizontal rules in tables

\usepackage{xcolor} % Required for specifying colors by name
\definecolor{ocre}{RGB}{243,102,25} % Define the orange color used for highlighting throughout the book

%----------------------------------------------------------------------------------------
%	MARGINS
%----------------------------------------------------------------------------------------

\usepackage{geometry} % Required for adjusting page dimensions and margins

\geometry{
	paper=a4paper, % Paper size, change to letterpaper for US letter size
	top=3cm, % Top margin
	bottom=3cm, % Bottom margin
	left=3cm, % Left margin
	right=3cm, % Right margin
	headheight=14pt, % Header height
	footskip=1.4cm, % Space from the bottom margin to the baseline of the footer
	headsep=10pt, % Space from the top margin to the baseline of the header
	%showframe, % Uncomment to show how the type block is set on the page
}

%----------------------------------------------------------------------------------------
%	FONTS
%----------------------------------------------------------------------------------------

\usepackage{avant} % Use the Avantgarde font for headings
%\usepackage{times} % Use the Times font for headings
\usepackage{mathptmx} % Use the Adobe Times Roman as the default text font together with math symbols from the Sym­bol, Chancery and Com­puter Modern fonts

\usepackage{microtype} % Slightly tweak font spacing for aesthetics
\usepackage[utf8]{inputenc} % Required for including letters with accents
\usepackage[T1]{fontenc} % Use 8-bit encoding that has 256 glyphs

%----------------------------------------------------------------------------------------
%	BIBLIOGRAPHY AND INDEX
%----------------------------------------------------------------------------------------

\usepackage[style=numeric,citestyle=numeric,sorting=nyt,sortcites=true,autopunct=true,babel=hyphen,hyperref=true,abbreviate=false,backref=true,backend=biber]{biblatex}
\addbibresource{bibliography.bib} % BibTeX bibliography file
\defbibheading{bibempty}{}

\usepackage{calc} % For simpler calculation - used for spacing the index letter headings correctly
\usepackage{makeidx} % Required to make an index
\makeindex % Tells LaTeX to create the files required for indexing

%----------------------------------------------------------------------------------------
%	MAIN TABLE OF CONTENTS
%----------------------------------------------------------------------------------------

\usepackage{titletoc} % Required for manipulating the table of contents

\contentsmargin{0cm} % Removes the default margin

% Part text styling (this is mostly taken care of in the PART HEADINGS section of this file)
\titlecontents{part}
	[0cm] % Left indentation
	{\addvspace{20pt}\bfseries} % Spacing and font options for parts
	{}
	{}
	{}

% Chapter text styling
\titlecontents{chapter}
	[1.25cm] % Left indentation
	{\addvspace{12pt}\large\sffamily\bfseries} % Spacing and font options for chapters
	{\color{ocre!60}\contentslabel[\Large\thecontentslabel]{1.25cm}\color{ocre}} % Formatting of numbered sections of this type
	{\color{ocre}} % Formatting of numberless sections of this type
	{\color{ocre!60}\normalsize\;\titlerule*[.5pc]{.}\;\thecontentspage} % Formatting of the filler to the right of the heading and the page number

% Section text styling
\titlecontents{section}
	[1.25cm] % Left indentation
	{\addvspace{3pt}\sffamily\bfseries} % Spacing and font options for sections
	{\contentslabel[\thecontentslabel]{1.25cm}} % Formatting of numbered sections of this type
	{} % Formatting of numberless sections of this type
	{\hfill\color{black}\thecontentspage} % Formatting of the filler to the right of the heading and the page number

% Subsection text styling
\titlecontents{subsection}
	[1.25cm] % Left indentation
	{\addvspace{1pt}\sffamily\small} % Spacing and font options for subsections
	{\contentslabel[\thecontentslabel]{1.25cm}} % Formatting of numbered sections of this type
	{} % Formatting of numberless sections of this type
	{\ \titlerule*[.5pc]{.}\;\thecontentspage} % Formatting of the filler to the right of the heading and the page number

% Figure text styling
\titlecontents{figure}
	[1.25cm] % Left indentation
	{\addvspace{1pt}\sffamily\small} % Spacing and font options for figures
	{\thecontentslabel\hspace*{1em}} % Formatting of numbered sections of this type
	{} % Formatting of numberless sections of this type
	{\ \titlerule*[.5pc]{.}\;\thecontentspage} % Formatting of the filler to the right of the heading and the page number

% Table text styling
\titlecontents{table}
	[1.25cm] % Left indentation
	{\addvspace{1pt}\sffamily\small} % Spacing and font options for tables
	{\thecontentslabel\hspace*{1em}} % Formatting of numbered sections of this type
	{} % Formatting of numberless sections of this type
	{\ \titlerule*[.5pc]{.}\;\thecontentspage} % Formatting of the filler to the right of the heading and the page number

%----------------------------------------------------------------------------------------
%	MINI TABLE OF CONTENTS IN PART HEADS
%----------------------------------------------------------------------------------------

% Chapter text styling
\titlecontents{lchapter}
	[0em] % Left indentation
	{\addvspace{15pt}\large\sffamily\bfseries} % Spacing and font options for chapters
	{\color{ocre}\contentslabel[\Large\thecontentslabel]{1.25cm}\color{ocre}} % Chapter number
	{}  
	{\color{ocre}\normalsize\sffamily\bfseries\;\titlerule*[.5pc]{.}\;\thecontentspage} % Page number

% Section text styling
\titlecontents{lsection}
	[0em] % Left indentation
	{\sffamily\small} % Spacing and font options for sections
	{\contentslabel[\thecontentslabel]{1.25cm}} % Section number
	{}
	{}

% Subsection text styling (note these aren't shown by default, display them by searchings this file for tocdepth and reading the commented text)
\titlecontents{lsubsection}
	[.5em] % Left indentation
	{\sffamily\footnotesize} % Spacing and font options for subsections
	{\contentslabel[\thecontentslabel]{1.25cm}}
	{}
	{}

%----------------------------------------------------------------------------------------
%	HEADERS AND FOOTERS
%----------------------------------------------------------------------------------------

\usepackage{fancyhdr} % Required for header and footer configuration

\pagestyle{fancy} % Enable the custom headers and footers

\renewcommand{\chaptermark}[1]{\markboth{\sffamily\normalsize\bfseries\chaptername\ \thechapter.\ #1}{}} % Styling for the current chapter in the header
\renewcommand{\sectionmark}[1]{\markright{\sffamily\normalsize\thesection\hspace{5pt}#1}{}} % Styling for the current section in the header

\fancyhf{} % Clear default headers and footers
\fancyhead[LE,RO]{\sffamily\normalsize\thepage} % Styling for the page number in the header
\fancyhead[LO]{\rightmark} % Print the nearest section name on the left side of odd pages
\fancyhead[RE]{\leftmark} % Print the current chapter name on the right side of even pages
%\fancyfoot[C]{\thepage} % Uncomment to include a footer

\renewcommand{\headrulewidth}{0.5pt} % Thickness of the rule under the header

\fancypagestyle{plain}{% Style for when a plain pagestyle is specified
	\fancyhead{}\renewcommand{\headrulewidth}{0pt}%
}

% Removes the header from odd empty pages at the end of chapters
\makeatletter
\renewcommand{\cleardoublepage}{
\clearpage\ifodd\c@page\else
\hbox{}
\vspace*{\fill}
\thispagestyle{empty}
\newpage
\fi}

%----------------------------------------------------------------------------------------
%	THEOREM STYLES
%----------------------------------------------------------------------------------------

\usepackage{amsmath,amsfonts,amssymb,amsthm} % For math equations, theorems, symbols, etc

\newcommand{\intoo}[2]{\mathopen{]}#1\,;#2\mathclose{[}}
\newcommand{\ud}{\mathop{\mathrm{{}d}}\mathopen{}}
\newcommand{\intff}[2]{\mathopen{[}#1\,;#2\mathclose{]}}
\renewcommand{\qedsymbol}{$\blacksquare$}
\newtheorem{notation}{Notation}[chapter]

% Boxed/framed environments
\newtheoremstyle{ocrenumbox}% Theorem style name
{0pt}% Space above
{0pt}% Space below
{\normalfont}% Body font
{}% Indent amount
{\small\bf\sffamily\color{ocre}}% Theorem head font
{\;}% Punctuation after theorem head
{0.25em}% Space after theorem head
{\small\sffamily\color{ocre}\thmname{#1}\nobreakspace\thmnumber{\@ifnotempty{#1}{}\@upn{#2}}% Theorem text (e.g. Theorem 2.1)
\thmnote{\nobreakspace\the\thm@notefont\sffamily\bfseries\color{black}---\nobreakspace#3.}} % Optional theorem note

\newtheoremstyle{blacknumex}% Theorem style name
{5pt}% Space above
{5pt}% Space below
{\normalfont}% Body font
{} % Indent amount
{\small\bf\sffamily}% Theorem head font
{\;}% Punctuation after theorem head
{0.25em}% Space after theorem head
{\small\sffamily{\tiny\ensuremath{\blacksquare}}\nobreakspace\thmname{#1}\nobreakspace\thmnumber{\@ifnotempty{#1}{}\@upn{#2}}% Theorem text (e.g. Theorem 2.1)
\thmnote{\nobreakspace\the\thm@notefont\sffamily\bfseries---\nobreakspace#3.}}% Optional theorem note

\newtheoremstyle{blacknumbox} % Theorem style name
{0pt}% Space above
{0pt}% Space below
{\normalfont}% Body font
{}% Indent amount
{\small\bf\sffamily}% Theorem head font
{\;}% Punctuation after theorem head
{0.25em}% Space after theorem head
{\small\sffamily\thmname{#1}\nobreakspace\thmnumber{\@ifnotempty{#1}{}\@upn{#2}}% Theorem text (e.g. Theorem 2.1)
\thmnote{\nobreakspace\the\thm@notefont\sffamily\bfseries---\nobreakspace#3.}}% Optional theorem note

% Non-boxed/non-framed environments
\newtheoremstyle{ocrenum}% Theorem style name
{5pt}% Space above
{5pt}% Space below
{\normalfont}% Body font
{}% Indent amount
{\small\bf\sffamily\color{ocre}}% Theorem head font
{\;}% Punctuation after theorem head
{0.25em}% Space after theorem head
{\small\sffamily\color{ocre}\thmname{#1}\nobreakspace\thmnumber{\@ifnotempty{#1}{}\@upn{#2}}% Theorem text (e.g. Theorem 2.1)
\thmnote{\nobreakspace\the\thm@notefont\sffamily\bfseries\color{black}---\nobreakspace#3.}} % Optional theorem note
\makeatother

% Defines the theorem text style for each type of theorem to one of the three styles above
\newcounter{dummy} 
\numberwithin{dummy}{section}
\theoremstyle{ocrenumbox}
\newtheorem{theoremeT}[dummy]{Theorem}
\newtheorem{problem}{Problem}[chapter]
\newtheorem{exerciseT}{Exercise}[chapter]
\theoremstyle{blacknumex}
\newtheorem{exampleT}{Example}[chapter]
\theoremstyle{blacknumbox}
\newtheorem{vocabulary}{Vocabulary}[chapter]
\newtheorem{definitionT}{Definition}[section]
\newtheorem{corollaryT}[dummy]{Corollary}
\theoremstyle{ocrenum}
\newtheorem{proposition}[dummy]{Proposition}

%----------------------------------------------------------------------------------------
%	DEFINITION OF COLORED BOXES
%----------------------------------------------------------------------------------------

\RequirePackage[framemethod=default]{mdframed} % Required for creating the theorem, definition, exercise and corollary boxes

% Theorem box
\newmdenv[skipabove=7pt,
skipbelow=7pt,
backgroundcolor=black!5,
linecolor=ocre,
innerleftmargin=5pt,
innerrightmargin=5pt,
innertopmargin=5pt,
leftmargin=0cm,
rightmargin=0cm,
innerbottommargin=5pt]{tBox}

% Exercise box	  
\newmdenv[skipabove=7pt,
skipbelow=7pt,
rightline=false,
leftline=true,
topline=false,
bottomline=false,
backgroundcolor=ocre!10,
linecolor=ocre,
innerleftmargin=5pt,
innerrightmargin=5pt,
innertopmargin=5pt,
innerbottommargin=5pt,
leftmargin=0cm,
rightmargin=0cm,
linewidth=4pt]{eBox}	

% Definition box
\newmdenv[skipabove=7pt,
skipbelow=7pt,
rightline=false,
leftline=true,
topline=false,
bottomline=false,
linecolor=ocre,
innerleftmargin=5pt,
innerrightmargin=5pt,
innertopmargin=0pt,
leftmargin=0cm,
rightmargin=0cm,
linewidth=4pt,
innerbottommargin=0pt]{dBox}	

% Corollary box
\newmdenv[skipabove=7pt,
skipbelow=7pt,
rightline=false,
leftline=true,
topline=false,
bottomline=false,
linecolor=gray,
backgroundcolor=black!5,
innerleftmargin=5pt,
innerrightmargin=5pt,
innertopmargin=5pt,
leftmargin=0cm,
rightmargin=0cm,
linewidth=4pt,
innerbottommargin=5pt]{cBox}

% Creates an environment for each type of theorem and assigns it a theorem text style from the "Theorem Styles" section above and a colored box from above
\newenvironment{theorem}{\begin{tBox}\begin{theoremeT}}{\end{theoremeT}\end{tBox}}
\newenvironment{exercise}{\begin{eBox}\begin{exerciseT}}{\hfill{\color{ocre}\tiny\ensuremath{\blacksquare}}\end{exerciseT}\end{eBox}}				  
\newenvironment{definition}{\begin{dBox}\begin{definitionT}}{\end{definitionT}\end{dBox}}	
\newenvironment{example}{\begin{exampleT}}{\hfill{\tiny\ensuremath{\blacksquare}}\end{exampleT}}		
\newenvironment{corollary}{\begin{cBox}\begin{corollaryT}}{\end{corollaryT}\end{cBox}}	

%----------------------------------------------------------------------------------------
%	REMARK ENVIRONMENT
%----------------------------------------------------------------------------------------

\newenvironment{remark}{\par\vspace{10pt}\small % Vertical white space above the remark and smaller font size
\begin{list}{}{
\leftmargin=35pt % Indentation on the left
\rightmargin=25pt}\item\ignorespaces % Indentation on the right
\makebox[-2.5pt]{\begin{tikzpicture}[overlay]
\node[draw=ocre!60,line width=1pt,circle,fill=ocre!25,font=\sffamily\bfseries,inner sep=2pt,outer sep=0pt] at (-15pt,0pt){\textcolor{ocre}{R}};\end{tikzpicture}} % Orange R in a circle
\advance\baselineskip -1pt}{\end{list}\vskip5pt} % Tighter line spacing and white space after remark

%----------------------------------------------------------------------------------------
%	SECTION NUMBERING IN THE MARGIN
%----------------------------------------------------------------------------------------

\makeatletter
\renewcommand{\@seccntformat}[1]{\llap{\textcolor{ocre}{\csname the#1\endcsname}\hspace{1em}}}                    
\renewcommand{\section}{\@startsection{section}{1}{\z@}
{-4ex \@plus -1ex \@minus -.4ex}
{1ex \@plus.2ex }
{\normalfont\large\sffamily\bfseries}}
\renewcommand{\subsection}{\@startsection {subsection}{2}{\z@}
{-3ex \@plus -0.1ex \@minus -.4ex}
{0.5ex \@plus.2ex }
{\normalfont\sffamily\bfseries}}
\renewcommand{\subsubsection}{\@startsection {subsubsection}{3}{\z@}
{-2ex \@plus -0.1ex \@minus -.2ex}
{.2ex \@plus.2ex }
{\normalfont\small\sffamily\bfseries}}                        
\renewcommand\paragraph{\@startsection{paragraph}{4}{\z@}
{-2ex \@plus-.2ex \@minus .2ex}
{.1ex}
{\normalfont\small\sffamily\bfseries}}

%----------------------------------------------------------------------------------------
%	PART HEADINGS
%----------------------------------------------------------------------------------------

% Numbered part in the table of contents
\newcommand{\@mypartnumtocformat}[2]{%
	\setlength\fboxsep{0pt}%
	\noindent\colorbox{ocre!20}{\strut\parbox[c][.7cm]{\ecart}{\color{ocre!70}\Large\sffamily\bfseries\centering#1}}\hskip\esp\colorbox{ocre!40}{\strut\parbox[c][.7cm]{\linewidth-\ecart-\esp}{\Large\sffamily\centering#2}}%
}

% Unnumbered part in the table of contents
\newcommand{\@myparttocformat}[1]{%
	\setlength\fboxsep{0pt}%
	\noindent\colorbox{ocre!40}{\strut\parbox[c][.7cm]{\linewidth}{\Large\sffamily\centering#1}}%
}

\newlength\esp
\setlength\esp{4pt}
\newlength\ecart
\setlength\ecart{1.2cm-\esp}
\newcommand{\thepartimage}{}%
\newcommand{\partimage}[1]{\renewcommand{\thepartimage}{#1}}%
\def\@part[#1]#2{%
\ifnum \c@secnumdepth >-2\relax%
\refstepcounter{part}%
\addcontentsline{toc}{part}{\texorpdfstring{\protect\@mypartnumtocformat{\thepart}{#1}}{\partname~\thepart\ ---\ #1}}
\else%
\addcontentsline{toc}{part}{\texorpdfstring{\protect\@myparttocformat{#1}}{#1}}%
\fi%
\startcontents%
\markboth{}{}%
{\thispagestyle{empty}%
\begin{tikzpicture}[remember picture,overlay]%
\node at (current page.north west){\begin{tikzpicture}[remember picture,overlay]%	
\fill[ocre!20](0cm,0cm) rectangle (\paperwidth,-\paperheight);
\node[anchor=north] at (4cm,-3.25cm){\color{ocre!40}\fontsize{220}{100}\sffamily\bfseries\thepart}; 
\node[anchor=south east] at (\paperwidth-1cm,-\paperheight+1cm){\parbox[t][][t]{8.5cm}{
\printcontents{l}{0}{\setcounter{tocdepth}{1}}% The depth to which the Part mini table of contents displays headings; 0 for chapters only, 1 for chapters and sections and 2 for chapters, sections and subsections
}};
\node[anchor=north east] at (\paperwidth-1.5cm,-3.25cm){\parbox[t][][t]{15cm}{\strut\raggedleft\color{white}\fontsize{30}{30}\sffamily\bfseries#2}};
\end{tikzpicture}};
\end{tikzpicture}}%
\@endpart}
\def\@spart#1{%
\startcontents%
\phantomsection
{\thispagestyle{empty}%
\begin{tikzpicture}[remember picture,overlay]%
\node at (current page.north west){\begin{tikzpicture}[remember picture,overlay]%	
\fill[ocre!20](0cm,0cm) rectangle (\paperwidth,-\paperheight);
\node[anchor=north east] at (\paperwidth-1.5cm,-3.25cm){\parbox[t][][t]{15cm}{\strut\raggedleft\color{white}\fontsize{30}{30}\sffamily\bfseries#1}};
\end{tikzpicture}};
\end{tikzpicture}}
\addcontentsline{toc}{part}{\texorpdfstring{%
\setlength\fboxsep{0pt}%
\noindent\protect\colorbox{ocre!40}{\strut\protect\parbox[c][.7cm]{\linewidth}{\Large\sffamily\protect\centering #1\quad\mbox{}}}}{#1}}%
\@endpart}
\def\@endpart{\vfil\newpage
\if@twoside
\if@openright
\null
\thispagestyle{empty}%
\newpage
\fi
\fi
\if@tempswa
\twocolumn
\fi}

%----------------------------------------------------------------------------------------
%	CHAPTER HEADINGS
%----------------------------------------------------------------------------------------

% A switch to conditionally include a picture, implemented by Christian Hupfer
\newif\ifusechapterimage
\usechapterimagetrue
\newcommand{\thechapterimage}{}%
\newcommand{\chapterimage}[1]{\ifusechapterimage\renewcommand{\thechapterimage}{#1}\fi}%
\newcommand{\autodot}{.}
\def\@makechapterhead#1{%
{\parindent \z@ \raggedright \normalfont
\ifnum \c@secnumdepth >\m@ne
\if@mainmatter
\begin{tikzpicture}[remember picture,overlay]
\node at (current page.north west)
{\begin{tikzpicture}[remember picture,overlay]
\node[anchor=north west,inner sep=0pt] at (0,0) {\ifusechapterimage\includegraphics[width=\paperwidth]{\thechapterimage}\fi};
\draw[anchor=west] (\Gm@lmargin,-9cm) node [line width=2pt,rounded corners=15pt,draw=ocre,fill=white,fill opacity=0.5,inner sep=15pt]{\strut\makebox[22cm]{}};
\draw[anchor=west] (\Gm@lmargin+.3cm,-9cm) node {\huge\sffamily\bfseries\color{black}\thechapter\autodot~#1\strut};
\end{tikzpicture}};
\end{tikzpicture}
\else
\begin{tikzpicture}[remember picture,overlay]
\node at (current page.north west)
{\begin{tikzpicture}[remember picture,overlay]
\node[anchor=north west,inner sep=0pt] at (0,0) {\ifusechapterimage\includegraphics[width=\paperwidth]{\thechapterimage}\fi};
\draw[anchor=west] (\Gm@lmargin,-9cm) node [line width=2pt,rounded corners=15pt,draw=ocre,fill=white,fill opacity=0.5,inner sep=15pt]{\strut\makebox[22cm]{}};
\draw[anchor=west] (\Gm@lmargin+.3cm,-9cm) node {\huge\sffamily\bfseries\color{black}#1\strut};
\end{tikzpicture}};
\end{tikzpicture}
\fi\fi\par\vspace*{270\p@}}}

%-------------------------------------------

\def\@makeschapterhead#1{%
\begin{tikzpicture}[remember picture,overlay]
\node at (current page.north west)
{\begin{tikzpicture}[remember picture,overlay]
\node[anchor=north west,inner sep=0pt] at (0,0) {\ifusechapterimage\includegraphics[width=\paperwidth]{\thechapterimage}\fi};
\draw[anchor=west] (\Gm@lmargin,-9cm) node [line width=2pt,rounded corners=15pt,draw=ocre,fill=white,fill opacity=0.5,inner sep=15pt]{\strut\makebox[22cm]{}};
\draw[anchor=west] (\Gm@lmargin+.3cm,-9cm) node {\huge\sffamily\bfseries\color{black}#1\strut};
\end{tikzpicture}};
\end{tikzpicture}
\par\vspace*{270\p@}}
\makeatother

%----------------------------------------------------------------------------------------
%	LINKS
%----------------------------------------------------------------------------------------

\usepackage{hyperref}
\hypersetup{hidelinks,backref=true,pagebackref=true,hyperindex=true,colorlinks=false,breaklinks=true,urlcolor=ocre,bookmarks=true,bookmarksopen=false}

\usepackage{bookmark}
\bookmarksetup{
open,
numbered,
addtohook={%
\ifnum\bookmarkget{level}=0 % chapter
\bookmarksetup{bold}%
\fi
\ifnum\bookmarkget{level}=-1 % part
\bookmarksetup{color=ocre,bold}%
\fi
}
}
 % Insert the commands.tex file which contains the majority of the structure behind the template

%\hypersetup{pdftitle={Title},pdfauthor={Author}} % Uncomment and fill out to include PDF metadata for the author and title of the book

%----------------------------------------------------------------------------------------

\begin{document}

%----------------------------------------------------------------------------------------
%	TITLE PAGE
%----------------------------------------------------------------------------------------

\begingroup
\thispagestyle{empty} % Suppress headers and footers on the title page
\begin{tikzpicture}[remember picture,overlay]
\node[inner sep=0pt] (background) at (current page.center) {\includegraphics[width=\paperwidth]{background.pdf}};
\draw (current page.center) node [fill=ocre!30!white,fill opacity=0.6,text opacity=1,inner sep=1cm]{\Huge\centering\bfseries\sffamily\parbox[c][][t]{\paperwidth}{\centering The Search for a Title\\[15pt] % Book title
{\Large A Profound Subtitle}\\[20pt] % Subtitle
{\huge Dr. John Smith}}}; % Author name
\end{tikzpicture}
\vfill
\endgroup

%----------------------------------------------------------------------------------------
%	COPYRIGHT PAGE
%----------------------------------------------------------------------------------------

\newpage
~\vfill
\thispagestyle{empty}

\noindent Copyright \copyright\ 2019 John Smith\\ % Copyright notice

\noindent \textsc{Published by Publisher}\\ % Publisher

\noindent \textsc{book-website.com}\\ % URL

\noindent Licensed under the Creative Commons Attribution-NonCommercial 3.0 Unported License (the ``License''). You may not use this file except in compliance with the License. You may obtain a copy of the License at \url{http://creativecommons.org/licenses/by-nc/3.0}. Unless required by applicable law or agreed to in writing, software distributed under the License is distributed on an \textsc{``as is'' basis, without warranties or conditions of any kind}, either express or implied. See the License for the specific language governing permissions and limitations under the License.\\ % License information, replace this with your own license (if any)

\noindent \textit{First printing, March 2019} % Printing/edition date

%----------------------------------------------------------------------------------------
%	TABLE OF CONTENTS
%----------------------------------------------------------------------------------------

%\usechapterimagefalse % If you don't want to include a chapter image, use this to toggle images off - it can be enabled later with \usechapterimagetrue

\chapterimage{chapter_head_1.pdf} % Table of contents heading image

\pagestyle{empty} % Disable headers and footers for the following pages

\tableofcontents % Print the table of contents itself

\cleardoublepage % Forces the first chapter to start on an odd page so it's on the right side of the book

\pagestyle{fancy} % Enable headers and footers again

%----------------------------------------------------------------------------------------
%	PART
%----------------------------------------------------------------------------------------

\part{Part One}

%----------------------------------------------------------------------------------------
%	CHAPTER 1
%----------------------------------------------------------------------------------------

\chapterimage{chapter_head_2.pdf} % Chapter heading image

\chapter{Text Chapter}
    
In this chapter, we will give the important results for the proof of the 
theorem of existence of optimal control pairs. In addition to presenting some 
problems related to these optimal pairs \cite{Yong2015}.

In what follows, we denote $\mathbb{R}_{+}=[0,\infty)$. For any initial pair 
$(t,x)\in \mathbb{R}_{+}\times \mathbb{R}^n$, consider the following ordinary 
differential equation:

\begin{equation}\label{eq1.10}
    \left\{ 
        \begin{array}{l}
            \dot{X}(s)=f(s,u(s),X(s)\,\,s\in [t,\infty), \\
            X(t)=x,\\
        \end{array}
        \right.
\end{equation}

Where $f:\mathbb{R}_{+}\times U\times \mathbb{R}^n\rightarrow \mathbb{R}^n$ is 
a given map. In the above, $X(\cdot)$ is called the state trajectory, taking 
values in $\mathbb{R}^n$ and $u(\cdot)$ is called the control, taking values in 
some metric space $U$. We call (\ref{eq1.10}) a control system.

For any $0\leq t<T<\infty$, we define the following:

$$\mathcal{U}[t,T]=\{u:[t,T]\rightarrow U | u(\cdot)\,\mbox{is measurable} \},$$

and 


$$
    \mathcal{U}[t,\infty)=
        \{u:[t,\infty)\rightarrow U | u(\cdot)\,
        \mbox{is  measurable} \}.
$$

Any $u(\cdot)\in \mathcal{U}[t,T]$, respectively 
$
    u(\cdot)\in 
    \mathcal{U}[t,\infty)
$, is called a feasible control on $[t,T]$, respectively $[t,\infty)$. 

Under proper conditions, for any initial pair $(t,x)$, and feasible control 
$u(\cdot)$, (\ref{eq1.10}) admits a unique solution 
$X(\cdot)=X(\cdot;t,x,u(\cdot))$ deined on $[t,\infty)$. Clearly different 
choices of $U(\cdot)$ will result in different state trajectories $X(\cdot)$. 
We reer to $(u(\cdot),X(\cdot))$ as a state-control pair of the control system 
(\ref{eq1.10}).

Next, we recall that $2^{\mathbb{R}^n}$ is the set of all subsets of 
$\mathbb{R}^n$. Any map $M:\mathbb{R}_{+}\rightarrow 2^{\mathbb{R}^n}$ is 
called a moving target in $\mathbb{R}^n$ i for any $t\in \mathbb{R}_{+}$, 
$M(t)$ is a measurable set in $\mathbb{R}^n$. WE allow $M(t)$ to be empty for 
some or all $t$, which will give us some flexibility below. In most situations, 
for any $t\in\mathbb{R}_{+}$, $M(t)$ is assumed to be closed or open.

Problem (C). Let $M(\cdot)$ be a moving target set in $\mathbb{R}^n$. For given 
$(t,x)\in \mathbb{R}_{+}\times \mathbb{R}^n$, ind a control $u(\cdot)\in 
\mathcal{U}[t,\infty)$ such that for some $\tau \geq t$,

$$
    X(\tau;t,x,u(\cdot))\in M(\tau).
$$

The above is called a controllability problem for the system (\ref{eq1.10}) 
with the moving target set $M(\cdot)$. For a moving set $M(\cdot)$ in 
$\mathbb{R}^n$, and $T\in (0\infty)$, we define
%
$$
    \mathcal{U}^{M(\cdot)}_x[t,T]=
        \{u(\cdot)\in \mathcal{U}[t,\infty) | X(\tau;t,x,u(\cdot))
        \in  M(\tau),\,\mbox{for some}\, \tau\in[t,T]\},\, (t,x)\in [0,T]\times 
        \mathbb{R}^n,
$$
and
%
$$
    \mathcal{U}^{M(\cdot)}_x[t,\infty)=\bigcup_{T \geq t}
    \mathcal{U}^{M(\cdot)}_x[t,T],\,\,\mbox{for every}\, 
    (t,x)
    \in
    \mathbb{R}_{+} \times \mathbb{R}^n.
$$
%
Let $M\subseteq \mathbb{R}^n$ be non-empty and closed, $T\in(0,\infty)$, and let
%
$$
    M(t)=
        \emptyset I_{\mathbb{R}_{+}
        \backslash{T}}(t)
        + MI_{\{T\}}(t),\,\, t\in \mathbb{R}_{+}.
$$

Then, for any $(t,x)\in [0,T)\times \mathbb{R}^n$,

$$\mathcal{U}^{M(\cdot)}_x[t,T]=\{u(\cdot)\in \mathcal{U}[t,T] | X(T;t,x,u(\cdot))\in M \}\equiv \tilde{\mathcal{U}}^M_x [t,T].$$

And define the cost functional

$$J(t,x;u(\cdot))=\int_{t}^{T}g(s,u(s),X(s))ds+h(T,X(T))\equiv J^T(t,x,u(\cdot)).$$

For some maps $g(\cdot)$ and $h(\cdot)$. Now we introduce the following problem.

Problem $(OC)^T$. For given $(t,x)\in \mathbb{R}_{+}\times \mathbb{R}^n$ with $\tilde{\mathcal{U}}^M_x[t,T]\neq\emptyset$, find a $\bar{u}(\cdot)\in \tilde{\mathcal{U}}^M_x[t,T]$ such that

\begin{equation*}
J^T(t,x;\bar{u}(\cdot))=\inf_{u(\cdot)\in \tilde{\mathcal{U}}^M_x[t,T]} J^T(t,x;u(\cdot))\equiv V(t,x).
\end{equation*}

The above is called an optimal control problem, with a fixed terminal time and a terminal state constraint.
        \section{Background}

    In this section, we present some results which will be used in the 
    following sections. These results are presented mostly in \cite{Yong2015}.
    Let us introduce some spaces. For any 
    $0\leq t< T<\infty$ and $1\leq p <\infty$, define
    
    \begin{equation*}
        C([t,T];\mathbb{R}^n)=\{
            \varphi:[t,T]\rightarrow \mathbb{R}^n | 
            \varphi(\cdot)\,\mbox{is continuous}
        \},
    \end{equation*}

\begin{equation*}
\displaystyle L^{\infty}(t,T;\mathbb{R}^n)=\{\varphi:[t,T]\rightarrow \mathbb{R}^n | \varphi(\cdot)\,\mbox{measurable},\, \displaystyle \mbox{esssup}_{s\in[t,T]}|\varphi(s)|<\infty\},
\end{equation*}

\begin{equation*}
\displaystyle L^{p}(t,T;\mathbb{R}^n)=\{\varphi:[t,T]\rightarrow \mathbb{R}^n | \varphi(\cdot)\,\mbox{measurable},\, \int_{t}^{T}|\varphi(s)|^p ds<\infty\},
\end{equation*}

which are Banach space under the following norms, respectively, 

\begin{equation*}
||\varphi(\cdot)||_{C([t,T];\mathbb{R}^n)}=\sup_{s\in[t,T]} |\varphi(s)|,\,\mbox{for every}\, \varphi(\cdot) \in C([t,T];\mathbb{R}^n),
\end{equation*}

\begin{equation*}
||\varphi(\cdot)||_{L^{\infty}(t,T;\mathbb{R}^n)}=\mbox{esssup}_{s\in[t,T]} |\varphi(s)|,\,\mbox{for every}\, \varphi(\cdot) \in L^{\infty}(t,T;\mathbb{R}^n),
\end{equation*}

where $\mbox{esssup} f:=\inf\{M | \mu(\{x: f(x)>M\})=0\}$,

\begin{equation*}
||\varphi(\cdot)||_{L^{p}(t,T;\mathbb{R}^n)}=\left(\int_{t}^{T}|\varphi(s)|^p ds\right)^\frac{1}{p},\,\mbox{for every}\, \varphi(\cdot) \in L^{p}(t,T;\mathbb{R}^n).
\end{equation*}

We now present some standard results.

\begin{theorem}\label{BFT}[Banach fixed point theorem]
	Let $\mathbb{X}$ be a Banach space, and $S: \mathbb{X} \rightarrow \mathbb{X}$ be a map satisfying 
	
	\begin{equation}\label{eq1.15}
	||S(x)-S(y)||\leq \alpha ||x-y||,\,\mbox{for every}\, x,y \in \mathbb{X},
	\end{equation}
with $\alpha \in (0,1)$. There exists a unique $\bar{x} \in \mathbb{X}$ such that $S(\bar{x})=\bar{x}$.
\end{theorem}
\begin{proof}
	Let see that the map $S$ is continuous. Given $\epsilon>0$, and $||x-y||<\delta$ with $\delta=\frac{\epsilon}{\alpha}$, we have by \ref{eq1.15}
	
	$$||S(x)-S(y)||\leq \alpha ||x-y||<\epsilon.$$
	
	For every $x,y \in \mathbb{X}$. Now pick any $x_0 \in \mathbb{X}$, and define the sequence $x_k=S^k(x_0)$, $k\geq 1$. 
%using telescopic sequence
Then for any $k,l\geq 1$,
\begin{eqnarray*}
||x_{k+l}-x_{k}||&\leq& ||\sum_{i=k+1}^{k+l}(x_i-x_{i-1})||=||\sum_{i=k+1}^{k+l}(S^i(x_0)-S^{i-1}(x_0))||\\
&\leq& \sum_{i=k+1}^{k+l}||(S^i(x_0)-S^{i-1}(x_0))||\leq \sum_{i=k+1}^{k+l} \alpha^k||x_1-x_0||.
\end{eqnarray*}

Thus, $\{x_k\}_{k>0}$ is a Cauchy sequence. Consequently, there exists a unique $\bar{x} \in \mathbb{X}$ such that 

$$\lim_{k\rightarrow \infty}||x_k-\bar{x}||=0.$$

Then by continuity of $S$, we obtain

\begin{center}
	$\bar{x}=\lim\limits_{k\rightarrow \infty}=\lim\limits_{k \rightarrow \infty} S(x_{k-1})=S(\bar{x}).$
\end{center}

This means that $\bar{x}$ is a fixed point of $S$. Finally, if $\bar{x}$ and $\tilde{x}$ are two fixed point. Then

$$||\bar{x}-\tilde{x}||=||S(\bar{x})-S(\tilde{x})||\leq \alpha ||\bar{x}-\tilde{x}||.$$

Hence, $\bar{x}=\tilde{x}$, proving the uniqueness.	
\end{proof}

\begin{theorem}\label{AAT}[Arzela-Ascoli]
	Let $\mathcal{Z}\subseteq C([t,T];\mathbb{R}^n)$ be an infinite set which is uniformly bounded and equicontinuous, i.e.
	$$\sup_{\varphi(\cdot)\in \mathcal{Z}}||\varphi(\cdot)||_{C([t,T];\mathbb{R}^n)}<\infty,$$
	
	and for any $\epsilon>0$, there exists a $\delta>0$ such that 
	\begin{center}
		$|\varphi(t)-\varphi(s)|<\epsilon$, for every $|t-s|<\delta$, $\varphi(\cdot)\in \mathcal{Z}$.
	\end{center}
	
	Then there exists a sequence $\varphi_k(\cdot)\in \mathcal{Z}$ such that 
	
	$$\lim_{k\rightarrow \infty}|\varphi_k(\cdot)-\bar{\varphi}(\cdot)||_{C([t,T];\mathbb{R}^n)}=0,$$

for some $\bar{\varphi}(\cdot) \in C([t,T];\mathbb{R}^n)$.	
\end{theorem}

\begin{proof}
Let define $D:=\{t_k\}_{k\geq 1}$ be a dense set of $[t,T]$. For any $k\geq 1$, the set $\{\varphi(t_1)|\varphi(\cdot)\in \mathcal{Z}\}$ is bounded. Thus, there exists aa sequence denoted by $\{\varphi_{\sigma_1(i)}(t_1)\}$ converging some point in $\mathbb{R}^n$, denoted by $\bar{\varphi}(t_1)$. Next, the set $\{\varphi_{\sigma_1(i)}(t_2)\}$ is bounded.

Thus, we may let $\{\varphi_{\sigma_2(i)}(t_2)\}$ be a subsequence of $\{\varphi_{\sigma_1(i)}(t_2)\}$, which is convergent to some point in $\mathbb{R}^n$ denoted by $\bar{\varphi}(t_2)$. Continue this process, we obtain a function $\bar{\varphi}:D\rightarrow \mathbb{R}$. by letting

$$\bar{\varphi}(\cdot)=\varphi_{\sigma_i(i)}(\cdot),\, i\geq 1.$$

We have 

$$\lim_{i\rightarrow \infty} \bar{\varphi}_i(s)=\bar{\varphi}(s)\,\mbox{for every}\, s\in D.$$

By equi-continuity of the sequence $\{\varphi_k(\cdot)\}$, we see that for any $\epsilon>0$ there exits a $\delta=\delta(\epsilon)>0$, independent of $i\geq 1$ such that

\begin{equation}\label{eq1.16}
	|\bar{\varphi}_i(s_1)-\bar{\varphi}_i(s_2)|<\epsilon\, \mbox{for every}\, s_1,s_2\in D, |s_1-s_2|<\delta.
\end{equation}

Then letting $i\rightarrow \infty$, we obtain

$$|\bar{\varphi}(s_1)-\bar{\varphi}(s_2)|\leq \epsilon\,\, \mbox{for every}\, s_1,s_2\in D, |s_1-s_2|<\delta.$$

This means that $\bar{\varphi}:D\rightarrow \mathbb{R}^n$ is uniformly continuous on $D$. Consequently, we may extend $\bar{\varphi}(\cdot)$ on $\bar{D}=[t,T]$ which is still continuous. Finally, for any $\epsilon>0$, let $\delta>0$ be such that (\ref{eq1.16}) holds and by compactness of $[t,T]$, let $S_{\delta}=\{s_j,1\leq j\leq M\}\subseteq D$ with $M>1$ depending on $\epsilon>0$ such that

$$[t,T]\subseteq \bigcup^M_{j=1}(s_j-\delta,s_j+\delta).$$

Next, we may let $i_0>1$ such that


	$$
	    \bar{\varphi}_i(s_j)-\bar{\varphi}(s_j)|<\epsilon, 
	    \quad
	    i\geq i_0, 1\leq j\leq M
	$$.


Then for any $s\in [t,T]$, there is an $s_j\in S_{\delta}$ such that $|s-s_j|<\delta$. Consequently 

$$|\bar{\varphi}_i(s)-\bar{\varphi}(s)|\leq |\bar{\varphi}_i(s)-\bar{\varphi}_i(s_j)|+|\bar{\varphi}_i(s_j)-\bar{\varphi}(s_j)|+|\bar{\varphi}(s_j)-\bar{\varphi}(s)|\leq 3\epsilon.$$

This show that $\varphi_i(\cdot)$ converges to $\bar{\varphi} (\cdot)$ uniformly  in $s\in [t,T]$.
\end{proof}

\begin{theorem}\label{BST}[Banach-Saks]
Let $\varphi_k(\cdot)\in L^2(a,b;\mathbb{R}^n)$ be a sequence which is weakly convergent to $\bar{\varphi}(\cdot)\in L^2(a,b;\mathbb{R}^n)$, i.e.,

$$\lim_{k\rightarrow \infty} \int_{a}^{b} \langle\varphi_{k}(s)-\bar{\varphi}(s),\eta(s)\rangle ds,\,\mbox{for every}\,\, \eta\in L^2(a,b;\mathbb{R}^n).$$

Then there is a subsequence $\{\varphi_{k_j}(\cdot)\}$ such that

$$\lim_{k\rightarrow \infty}||\frac{1}{N}\sum_{j=1}^{N}\varphi_{k_j}(\cdot)-\bar{\varphi}(\cdot)||_{L^2(a,b;\mathbb{R}^n)}=0.$$
\end{theorem}
\begin{proof}
	Whitout loss of generality, first we consider that $\bar{\varphi}(\cdot)=0$. Let $k_1=1$. By the weak convergence of $\varphi_k(\cdot)$, we may find $k_1<k_2<\cdots <k_N$ such that
	
	$$|\int_{a}^{b}\langle\varphi_{k_i}(s),\varphi_{k_j}(s)\rangle ds|<\frac{1}{N}\,\, 1\leq i<j\leq N$$
	
	observe
	\begin{eqnarray*}
	||\frac{1}{N}\sum_{i=1}^{N}\varphi_{k_i}(\cdot)||^2_{L^2(a,b;\mathbb{R}^2)}&=&\frac{1}{N^2}\int_{a}^{b}|\sum_{i=1}^{N}\varphi_{k_i}(s)|^2 ds = \frac{1}{N^2}\int \sum_{i,j=1}^{N}\langle\varphi_{k_i}(s),\varphi_{k_j}(s)\rangle ds\\
	&=&\frac{1}{N^2}\sum_{i=1}^{N}||\varphi_{k_i}(\cdot)||^2_{L^2(a,b;\mathbb{R}^n)}+\frac{2}{N^2}\sum_{1\leq i<j\leq N}\int_{a}^{b}\langle\varphi_{k_i}(s),\varphi_{k_j}(s)\rangle ds\\
	&\leq& \frac{1}{N}\sup_{i\geq 1}||\varphi_{k_i}(\cdot)||^2_{L^2(a,b;\mathbb{R}^n)}+\frac{2}{N^3}\frac{N(N-1)}{2}\\
	&\leq& \frac{1}{N}\sup_{i\geq 1} ||\varphi_{k_i}(\cdot)||^2_{L^2(a,b;\mathbb{R}^n)}+\frac{1}{N}\rightarrow 0,
	\end{eqnarray*}
	
	when $N\rightarrow 0$. Now, consider that $\bar{\varphi}(\cdot)\neq 0$, thus $\varphi_{k}(\cdot)-\bar{\varphi}(\cdot)=0$, and we can apply the previous steps.
\end{proof}

\begin{lemma}
    \label{FL}[Filippov]
    Let $U$ be a complete separable metric space whose metric is denoted by 
    $d(\cdot,\cdot)$. Let $g:[0,T]\times U\rightarrow \mathbb{R}^n$ be a map 
    which is measurable in $t\in[0,T]$ and
    
    \begin{center}
        $|g(t,u)-g(t,v)|\leq \omega(d(u,v))$, for every $u,v\in U$, $t\in 
        [0,t],$
    \end{center}
%
     for some continuous and increasing function 
     $\omega:\mathbb{R}_{+}\rightarrow \mathbb{R}_{+}$ with $\omega(0)=0$, 
     called a modulus of continuity. Moreover,
%
    \begin{equation*}
        0\in g(t,U)\,\,\,a.e. \,\,t\in [0,t].
    \end{equation*}
    
    Then there exists a measurable map $u:[0,T]\rightarrow U$, such that 
    
    \begin{equation}\label{eq1.17}
        g(t,u(t))=0\,\, a.e.\,\, t\in [0,T].
    \end{equation}
\end{lemma}

\begin{proof}
Define $\bar{d}(u,v)=\frac{d(u,v)}{1+d(u,v)}<1$ for every $u,v\in U$. Then with this new metric $\bar{d}$, $U$ still complet and separable. Hence without loss of generality, we assume that the original metric $d(\cdot,\cdot)$ already satisfies $d(u,v)<1$ for every $u,v \in U$.

Next, we define

\begin{equation*}
\Gamma(t):=\{u\in U | g(t,u)=0\},\, t\in [0,t].
\end{equation*}	

We have that $\Gamma(t)\neq \emptyset$, because $0\in g(t,U)\,\,\,a.e. \,\,t\in [0,t]$. Let $U_0:=\{v_k | k\geq 1\}$ be a countable dense subset of $U$. We claim that for any $u \in U$ and $0\leq c < 1$,

\begin{equation}\label{eq1.18}
\{t\in[0,T] |\, d(u,\Gamma(t))\leq c\}=\bigcap^{\infty}_{i=1}\bigcup^{\infty}_{j=1}\{t\in[0,T] |\, d(u,v_j)\leq c+\frac{1}{i},|g(t,v_j)|\leq \frac{1}{i}\},
\end{equation}

where 

$$d(u,\Gamma (t)):=\inf_{v\in \Gamma(t)}d(u,v).$$

To show (\ref{eq1.18}), we note that $t\in [0,T]$, with $d(u,\Gamma(t))\leq c$ if and only if there exists a sequence $u_k \in \Gamma(t)$, i.e., $g(t,u_k)=0$, such that

$$d(u,u_k)\leq c+\frac{1}{k}.$$

Since $\bar{U}_0=U$, there exists a sequence $v_{j_k} \in U_0$ such that 

$$d(u_k,v_{j_k})<\frac{1}{k}.$$

Hence, by triangle inequality

$$d(u,v_{j_k})\leq d(u,u_k)+d(u_k,v_{j_k})\leq c+\frac{2}{k}.$$

Next, by the uniform continuity of $u\mapsto g(t,u)$, we have

$$|g(t,v_{j_k})|\leq |g(t,v_{j_k})-g(t,u_k)|\leq \omega(d(v_{j_k},u_k))\leq \omega\left(\frac{1}{k}\right).$$

Hence, one has

$$\left\{ \begin{array}{l}
	\lim_{k\rightarrow \infty} d(u,v_{j_k})\leq c, \\
    \lim_{k\rightarrow \infty} g(t,v_{j_k})=0. \\
\end{array}
\right.$$

Thus, $d(u,\Gamma(t))\leq c$ if and only if for any $i\geq 1$, there exists a $j\geq i$, such that

$$\left\{ \begin{array}{l}
 \bar{d}(u,v_{j})\leq c+\frac{1}{i}, \\
 |g(t,v_{j})|\leq \frac{1}{i}. \\
\end{array}
\right.$$

This prove (\ref{eq1.18}). Since the right-hand side of (\ref{eq1.18}) is measurable, so is the let-hand side. On the other hand,

$$\left\{ \begin{array}{l}
\{t\in[0,T] |\, d(u,\Gamma(t))\leq c\}=[0,T],\,\mbox{for every}\, c\geq 1, \\
\{t\in[0,T] |\, d(u,\Gamma(t))\leq c\}=\emptyset,\,\mbox{for every}\, c<0. \\
\end{array}
\right.$$

Hence, the function $t\mapsto d(u,\Gamma(t))$ is measurable. Now, we define

$$u_0(t):=v_1,\,\mbox{for every}\, t\in [0,T].$$

Clearly, $u_0(t)$ is measurable and 

$$d(u_0(t),\Gamma(t))<1\,\mbox{for every}\,t\in[0,T].$$

Suppose that we have defined $u_{k-1}(\cdot)$ such that 

\begin{equation}\label{eq1.19}
\left\{ \begin{array}{l}
d(u_{k-1}(t),\Gamma(t))\leq 2^{1-k},  \\
d(u_{k-1}(t),u_{k-2}(t))\leq 2^{2-k},\\
\end{array}
\right. t\in[0,T].
\end{equation}

We define the sets

$$\left\{ \begin{array}{l}
	C^k_i:=\{t\in[0,T] |\, d(v_i,\Gamma(t))<2^{-k}\}, \\
	D^k_i:=\{t\in[0,T] |\, d(v_i,u_{k-1}(t))< 2^{1-k}\}.\\
\end{array}
\right.$$

Since $t\mapsto d(v_i,\Gamma(t))$ is measurable, $C^k_i$ is measurable. Likewise, $D^k_i$ is also measurable. Set
$$A^k_i=C^k_i\cap D^k_i, \,k,i\geq 1.$$

Then $A^k_i$ is measurable as well. We claim that

\begin{equation}\label{eq1.20}
[0,T]=\bigcup^{\infty}_{i=1}A^k_i\,\mbox{for every}\,k\geq 1.
\end{equation}

In fact, for any $t\in[0,T]$, by (\ref{eq1.19}), there exists a $u\in\Gamma(t)$ such that 

$$d(u_{k-1}(t),u)<2^{1-k}.$$

By the density of $U_0$ in $U$, there exists an $i\geq 1$ such that

$$d(v_i,\Gamma(t))\leq\left\{ \begin{array}{l}
	d(v_i,u)< 2^{-k},  \\
	d(v_i,u_{k-1}(t))<2^{1-k},\\
\end{array}
\right.$$

which means $t\in A^k_i$, proving (\ref{eq1.20}). Now we define $u_k(\cdot):[0,T]\rightarrow U_0\subseteq U$ as follows:

$$u_k(t)=v_i,\,\mbox{for every}\, t\in A^k_{i}\backslash \bigcup^{\infty}_{j=1}A^k_j.$$

By $t\in C^k_i$, we have 

$$d(u_k(t),\Gamma(t))<2^{-k},$$

and by $t\in D^k_i$, we have

$$d(u_k(t),u_{k-1}(t))<2^{1-k}.$$

This completes the construction of the sequence $\{u_k(\cdot)\}$ inductively. Clearly, (\ref{eq1.19}) holds for every $k\geq 1$. This als implies that for each $t\in[0,T]$, $\{u_k(t)\}$ is Cauchy in $U$. By completeness of $U$m we obtain 

$$\lim_{k\rightarrow \infty} u_k(t)=u(t),\, t\in [0,t].$$

Of course, $u(\cdot)$ is measurable, and moreover, by the closeness of $\Gamma(t)$, we have

$$u(t)\in \Gamma(t),\,\mbox{for every}\, t\in [0,T].$$

This means (\ref{eq1.17}) holds.
\end{proof}

\begin{prop}\label{GIP}[Gronwall's Inequality]
	Let $\theta:[a,b]\rightarrow \mathbb{R}_{+}$ be continuous and satisfy
	$$\theta (s)\leq \alpha(s)+\int_a^s\beta(r)\theta (r) dr,\,\, s\in [a,b],$$
	
	for some $\alpha(\cdot),\beta (\cdot)\in L^1(a,b;\mathbb{R}_{+})$. Then
	
	\begin{equation}\label{eq1.30}
	\theta(s)\leq \alpha(s)+\int_{a}^{s}\alpha(\tau)\beta(\tau)e^{\int_{\tau}^{s}\beta(r)dr}d\tau,\,\, s\in[a,b].
	\end{equation}
	In particular, if $\alpha(\cdot)=\alpha$ is a constant, then
	
	\begin{equation}\label{eq1.31}
	\theta(s) \leq \alpha e^{\int_{a}^{s}\beta(r)dr},\,\, s[a,b]. 
	\end{equation}
	\end{prop}
\begin{proof}
	Let $\varphi(s)=\int_{a}^{s}\beta(r)\theta(r)dr.$, by the fundamental theorem of calculus, we have
		
	$$\dot{\varphi}(s)=\beta(s)\theta(s)\leq \beta(s)[\alpha(s)+\varphi(s)].$$
	
	This leads to 
	
	$$[\varphi(s)e^{-\int_{a}^{s}\beta(r)dr}]'\leq\alpha(s)\beta(s)e^{-\int_{a}^{s}\beta(r)dr}.$$
	
	Consequently,
	$$\varphi(s)e^{-\int_{a}^{s}\beta(r)dr}\leq \int_{a}^{s}\alpha(\tau)\beta(\tau)e^{-\int_{a}^{\tau}\beta(r)dr}d\tau.$$
	
		$$\varphi(s)\leq \int_{a}^{s}\alpha(\tau)\beta(\tau)e^{-\int_{a}^{\tau}\beta(r)dr}e^{\int_{a}^{s}\beta(r)dr}d\tau .$$
	
	rewriting 
	
	$$\varphi(s)\leq \int_{a}^{s}\alpha(\tau)\beta(\tau)e^{\int_{\tau}^{s}\beta(r)dr}d\tau .$$
	
	Hence,
	
	$$\theta(s)\leq \alpha(s)+\int_{a}^{s}\alpha(\tau)\beta(\tau)e^{\int_{\tau}^{s}\beta(r)dr}d\tau.$$
	
	(\ref{eq1.30}) holds. Now, consider $\alpha$ constant, then
	
	$$\theta(s)\leq \alpha+\int_{a}^{s}\alpha\beta(\tau)e^{\int_{\tau}^{s}\beta(r)dr}d\tau.$$
	
	By integration rules, if $u=\alpha,dv=\beta(\tau)e^{\int_{\tau}^{s}\beta(r)dr}d\tau$. In the other hand,
	
	$$\frac{d}{d\tau}e^{\int_{\tau}^{s}\beta(r)dr}=\frac{d}{d\tau}e^{-\int_{s}^{\tau}\beta(r)dr}=-\beta(\tau) e^{-\int_{s}^{\tau} \beta(r)dr}=-\beta(\tau)e^{\int_{\tau}^{s}\beta(r)dr}$$
	
	Then
	
	$$\theta(s)\leq \alpha-\alpha\int_{a}^{s}\frac{d}{d\tau}e^{\int_{\tau}^{s}\beta(r)dr}d\tau=\alpha-\alpha [e^{0}-e^{\int_{a}^{s}\beta(r)dr}].$$
	
	Therefore,
	
	$$\theta(s)\leq \alpha e^{\int_{a}^{s}\beta(r)dr},$$
	
	(\ref{eq1.31}) holds.
\end{proof}

%------------------------------------------------

\section{Citation}\index{Citation}

This statement requires citation \cite{article_key}; this one is more specific \cite[162]{book_key}.

%------------------------------------------------

\section{Lists}\index{Lists}

Lists are useful to present information in a concise and/or ordered way\footnote{Footnote example...}.

\subsection{Numbered List}\index{Lists!Numbered List}

\begin{enumerate}
\item The first item
\item The second item
\item The third item
\end{enumerate}

\subsection{Bullet Points}\index{Lists!Bullet Points}

\begin{itemize}
\item The first item
\item The second item
\item The third item
\end{itemize}

\subsection{Descriptions and Definitions}\index{Lists!Descriptions and Definitions}

\begin{description}
\item[Name] Description
\item[Word] Definition
\item[Comment] Elaboration
\end{description}

%----------------------------------------------------------------------------------------
%	CHAPTER 2
%----------------------------------------------------------------------------------------

\chapter{In-text Elements}

\section{Theorems}\index{Theorems}

This is an example of theorems.

\subsection{Several equations}\index{Theorems!Several Equations}
This is a theorem consisting of several equations.

\begin{theorem}[Name of the theorem]
In $E=\mathbb{R}^n$ all norms are equivalent. It has the properties:
\begin{align}
& \big| ||\mathbf{x}|| - ||\mathbf{y}|| \big|\leq || \mathbf{x}- \mathbf{y}||\\
&  ||\sum_{i=1}^n\mathbf{x}_i||\leq \sum_{i=1}^n||\mathbf{x}_i||\quad\text{where $n$ is a finite integer}
\end{align}
\end{theorem}

\subsection{Single Line}\index{Theorems!Single Line}
This is a theorem consisting of just one line.

\begin{theorem}
A set $\mathcal{D}(G)$ in dense in $L^2(G)$, $|\cdot|_0$. 
\end{theorem}

%------------------------------------------------

\section{Definitions}\index{Definitions}

This is an example of a definition. A definition could be mathematical or it could define a concept.

\begin{definition}[Definition name]
Given a vector space $E$, a norm on $E$ is an application, denoted $||\cdot||$, $E$ in $\mathbb{R}^+=[0,+\infty[$ such that:
\begin{align}
& ||\mathbf{x}||=0\ \Rightarrow\ \mathbf{x}=\mathbf{0}\\
& ||\lambda \mathbf{x}||=|\lambda|\cdot ||\mathbf{x}||\\
& ||\mathbf{x}+\mathbf{y}||\leq ||\mathbf{x}||+||\mathbf{y}||
\end{align}
\end{definition}

%------------------------------------------------

\section{Notations}\index{Notations}

\begin{notation}
Given an open subset $G$ of $\mathbb{R}^n$, the set of functions $\varphi$ are:
\begin{enumerate}
\item Bounded support $G$;
\item Infinitely differentiable;
\end{enumerate}
a vector space is denoted by $\mathcal{D}(G)$. 
\end{notation}

%------------------------------------------------

\section{Remarks}\index{Remarks}

This is an example of a remark.

\begin{remark}
The concepts presented here are now in conventional employment in mathematics. Vector spaces are taken over the field $\mathbb{K}=\mathbb{R}$, however, established properties are easily extended to $\mathbb{K}=\mathbb{C}$.
\end{remark}

%------------------------------------------------

\section{Corollaries}\index{Corollaries}

This is an example of a corollary.

\begin{corollary}[Corollary name]
The concepts presented here are now in conventional employment in mathematics. Vector spaces are taken over the field $\mathbb{K}=\mathbb{R}$, however, established properties are easily extended to $\mathbb{K}=\mathbb{C}$.
\end{corollary}

%------------------------------------------------

\section{Propositions}\index{Propositions}

This is an example of propositions.

\subsection{Several equations}\index{Propositions!Several Equations}

\begin{proposition}[Proposition name]
It has the properties:
\begin{align}
& \big| ||\mathbf{x}|| - ||\mathbf{y}|| \big|\leq || \mathbf{x}- \mathbf{y}||\\
&  ||\sum_{i=1}^n\mathbf{x}_i||\leq \sum_{i=1}^n||\mathbf{x}_i||\quad\text{where $n$ is a finite integer}
\end{align}
\end{proposition}

\subsection{Single Line}\index{Propositions!Single Line}

\begin{proposition} 
Let $f,g\in L^2(G)$; if $\forall \varphi\in\mathcal{D}(G)$, $(f,\varphi)_0=(g,\varphi)_0$ then $f = g$. 
\end{proposition}

%------------------------------------------------

\section{Examples}\index{Examples}

This is an example of examples.

\subsection{Equation and Text}\index{Examples!Equation and Text}

\begin{example}
Let $G=\{x\in\mathbb{R}^2:|x|<3\}$ and denoted by: $x^0=(1,1)$; consider the function:
\begin{equation}
f(x)=\left\{\begin{aligned} & \mathrm{e}^{|x|} & & \text{si $|x-x^0|\leq 1/2$}\\
& 0 & & \text{si $|x-x^0|> 1/2$}\end{aligned}\right.
\end{equation}
The function $f$ has bounded support, we can take $A=\{x\in\mathbb{R}^2:|x-x^0|\leq 1/2+\epsilon\}$ for all $\epsilon\in\intoo{0}{5/2-\sqrt{2}}$.
\end{example}

\subsection{Paragraph of Text}\index{Examples!Paragraph of Text}

\begin{example}[Example name]
\lipsum[2]
\end{example}

%------------------------------------------------

\section{Exercises}\index{Exercises}

This is an example of an exercise.

\begin{exercise}
This is a good place to ask a question to test learning progress or further cement ideas into students' minds.
\end{exercise}

%------------------------------------------------

\section{Problems}\index{Problems}

\begin{problem}
What is the average airspeed velocity of an unladen swallow?
\end{problem}

%------------------------------------------------

\section{Vocabulary}\index{Vocabulary}

Define a word to improve a students' vocabulary.

\begin{vocabulary}[Word]
Definition of word.
\end{vocabulary}

%----------------------------------------------------------------------------------------
%	PART
%----------------------------------------------------------------------------------------

\part{Part Two}

%----------------------------------------------------------------------------------------
%	CHAPTER 3
%----------------------------------------------------------------------------------------

\chapterimage{chapter_head_1.pdf} % Chapter heading image

\chapter{Presenting Information}

\section{Table}\index{Table}

\begin{table}[h]
\centering
\begin{tabular}{l l l}
\toprule
\textbf{Treatments} & \textbf{Response 1} & \textbf{Response 2}\\
\midrule
Treatment 1 & 0.0003262 & 0.562 \\
Treatment 2 & 0.0015681 & 0.910 \\
Treatment 3 & 0.0009271 & 0.296 \\
\bottomrule
\end{tabular}
\caption{Table caption}
\label{tab:example} % Unique label used for referencing the table in-text
%\addcontentsline{toc}{table}{Table \ref{tab:example}} % Uncomment to add the table to the table of contents
\end{table}

Referencing Table \ref{tab:example} in-text automatically.

%------------------------------------------------

\section{Figure}\index{Figure}

\begin{figure}[h]
\centering\includegraphics[scale=0.5]{placeholder.jpg}
\caption{Figure caption}
\label{fig:placeholder} % Unique label used for referencing the figure in-text
%\addcontentsline{toc}{figure}{Figure \ref{fig:placeholder}} % Uncomment to add the figure to the table of contents
\end{figure}

Referencing Figure \ref{fig:placeholder} in-text automatically.

%----------------------------------------------------------------------------------------
%	BIBLIOGRAPHY
%----------------------------------------------------------------------------------------

\chapter*{Bibliography}
\addcontentsline{toc}{chapter}{\textcolor{ocre}{Bibliography}} % Add a Bibliography heading to the table of contents

%------------------------------------------------

\section*{Articles}
\addcontentsline{toc}{section}{Articles}
\printbibliography[heading=bibempty,type=article]

%------------------------------------------------

\section*{Books}
\addcontentsline{toc}{section}{Books}
\printbibliography[heading=bibempty,type=book]

%----------------------------------------------------------------------------------------
%	INDEX
%----------------------------------------------------------------------------------------

\cleardoublepage % Make sure the index starts on an odd (right side) page
\phantomsection
\setlength{\columnsep}{0.75cm} % Space between the 2 columns of the index
\addcontentsline{toc}{chapter}{\textcolor{ocre}{Index}} % Add an Index heading to the table of contents
\printindex % Output the index

%----------------------------------------------------------------------------------------

\end{document}
